\sgaStableTarget{sga6.e08.expose}\section*{\underline{EXPOSÉ VIII}\\[3pt] \underline{\underline{THE RIEMANN--ROCH THEOREM}}}
\begin{center}
\underline{by Pierre BERTHELOT}
\end{center}

\sgaStableTarget{sga6.e08.sec.1}\subsection*{1. Complete-intersection morphisms}

\sgaStableTarget{sga6.e08.def.1-1}\textbf{Definition 1.1.} \uline{Let} @@M0@@ \uline{be a morphism of schemes}. \uline{One says that} @@M1@@ \uline{is a complete-intersection morphism if for every point} @@M2@@ \uline{of} @@M3@@, \uline{there exist an open neighbourhood} @@M4@@ \uline{of} @@M5@@, \uline{a smooth} @@M6@@-\uline{scheme} @@M7@@, \uline{such that the restriction of} @@M8@@ \uline{to} @@M9@@ \uline{factors as}:
@@M10@@
\uline{where} @@M11@@ \uline{is a regular immersion} (VII 1.4).

After replacing @@M12@@ by an open subset, one may suppose that @@M13@@ is a regular closed immersion.

In this definition one may replace @@M14@@ by a vector bundle @@M15@@. Indeed, since the notion is local, one may suppose that @@M16@@ is affine, and that its image in @@M17@@ is contained in an affine open subset @@M18@@ of @@M19@@. There then exist an integer @@M20@@ and a closed @@M21@@-immersion @@M22@@, hence a @@M23@@-immersion @@M24@@, which is regular by VII 1.10. Since the composite of two regular immersions is regular, one thus obtains a regular @@M25@@-immersion @@M26@@ factoring @@M27@@.

\sgaStableTarget{sga6.e08.prop.1-2}\textbf{Proposition 1.2.} \uline{Let} @@M28@@ \uline{be a complete-intersection morphism}. \uline{Then every} @@M29@@-\uline{immersion} @@M30@@ \uline{of} @@M31@@ \uline{into a smooth} @@M32@@-\uline{scheme is regular}.

We shall first prove the corollary:

\sgaStableTarget{sga6.e08.cor.1-3}\textbf{Corollary 1.3.} \uline{Consider the commutative diagram}:
@@M33@@
